\chapter*{Thư Gửi Phụ Huynh Và Giáo Viên}
\addcontentsline{toc}{chapter}{Thư Gửi Phụ Huynh Và Giáo Viên}
\markboth{Thư Gửi Phụ Huynh Và Giáo Viên}{Thư Gửi Phụ Huynh Và Giáo Viên}

\begin{truyen}{Gửi Những Người Đang Đồng Hành Với Trẻ}{To The Adults Walking Beside Young People}
\chuhoa{N}{ếu học sinh của hôm nay} làm người lớn mệt hơn trước, điều đó không chỉ vì các em yếu hơn hay bướng hơn. Các em đang lớn lên trong một thời đại nhiều chữ hơn, nhiều màn hình hơn, nhiều ví dụ thành công giả lập hơn, và cũng nhiều cơ hội tự biện hộ cho mình hơn. Một đứa trẻ mười bảy tuổi bây giờ có thể nghe hàng trăm lời khuyên trong một ngày, nhưng lại thiếu vài cuộc trò chuyện thật sự khiến nó chịu dừng lại để nghĩ.
\textit{(If students today exhaust adults more than before, it is not simply because they are weaker or more defiant. They are growing up in an age with more language, more screens, more simulated success stories, and more tools to justify themselves. A seventeen-year-old may hear hundreds of pieces of advice in a day, yet still lack a few real conversations that make them stop and think.)}

Quyển sách này không xin phụ huynh và giáo viên phải mềm mỏng vô điều kiện. Nó cũng không xin người lớn nhường hết cho tâm lý tuổi trẻ. Điều nó muốn là một cách cứng hơn nhưng khôn hơn: vẫn giữ chuẩn, nhưng bớt làm nhục; vẫn sửa sai, nhưng bớt nói theo kiểu đóng cửa đối thoại; vẫn kỳ vọng, nhưng đừng để kỳ vọng biến thành lời tuyên án lặp lại mỗi ngày.
\textit{(This book does not ask parents and teachers to become endlessly soft. Nor does it ask adults to surrender everything to youth psychology. It asks for something firmer yet wiser: keep standards, but humiliate less; correct mistakes, but stop speaking in ways that shut dialogue down; keep expectations, but do not let them become daily verdicts.)}
\end{truyen}

\section*{Ba Điều Người Lớn Có Thể Làm Ngay}

\begin{goigiaovien}
\textbf{Thứ nhất: đổi câu mở đầu.}

Thay vì bắt đầu bằng ``Em lúc nào cũng...'' hoặc ``Con chẳng bao giờ...'', hãy bắt đầu bằng một việc cụ thể vừa xảy ra. Câu mở đầu càng chính xác, cuộc nói chuyện càng ít biến thành chiến tranh tự ái.

\textbf{Thứ hai: phân biệt hành vi với con người.}

Một học sinh vô trách nhiệm trong một giai đoạn không có nghĩa là em ấy là người bỏ đi. Một lỗi đi học muộn, nói dối, trốn bài hay chống chế cần bị chỉnh, nhưng không nên biến thành nhãn dán lâu dài lên nhân phẩm.

\textbf{Thứ ba: để lại một cánh cửa quay về.}

Sau khi nhắc nhở hay xử lý kỷ luật, người lớn nên để học sinh thấy vẫn còn đường sửa. Một đứa trẻ tin rằng mình còn cơ hội sẽ dễ hợp tác hơn một đứa trẻ nghĩ mình đã bị xếp loại xong.
\end{goigiaovien}

\begin{baihoc}
Người lớn có quyền đặt giới hạn. Nhưng giới hạn mạnh nhất không phải giới hạn làm trẻ sợ; đó là giới hạn khiến trẻ hiểu rằng nếu bước qua, các em phải chịu trách nhiệm, nhưng vẫn còn cơ hội trưởng thành.
\textit{(Adults have the right to set limits. But the strongest limit is not the one that merely scares the young; it is the one that makes them understand that if they cross it, they must face responsibility, yet still have a chance to grow.)}
\end{baihoc}

\begin{ghinhoanh}
\textbf{A note to adults:}

Be firm without humiliation.

Leave a door open for change.
\end{ghinhoanh}

\begin{tuhocnhanh}
1. Trong cách nói chuyện với học sinh hoặc con cái, tôi đang lặp câu nào nhiều nhất mà ít hiệu quả nhất? \textit{(Which sentence do I repeat most often with students or children, and with the least effect?)}

2. Tôi đang giữ nguyên tắc, hay đang giữ sĩ diện của người lớn? \textit{(Am I protecting principles, or my adult pride?)}

3. Sau khi đọc quyển này, tôi có thể đổi ngay một cách nói nào để cuộc đối thoại bớt đóng lại? \textit{(After reading this book, what one phrase can I change immediately so the conversation shuts down less?)}
\end{tuhocnhanh}

\ngancach